% !TEX root = ../algorithms.tex

\section{Хеш-таблицы, ключи. Амортизированное и ожидаемое время работы. Время работы хеш-таблицы с открытой адресацией при идеализированном случайном равновероятном хешировании}

\begin{Definition}{Словарь (хеш-таблица)}{}
	\emph{Словарь} --- структура данных для хранения пар <<ключ--значение>> $(x \in U,\; v)$, поддерживающая операции:
	\begin{itemize}
		\item $\mathrm{insert}(x,\,v)$ --- вставить значение $v$ по ключу $x$;
		\item $\mathrm{delete}(x)$ --- удалить запись с ключом $x$;
		\item $v = \mathrm{query}(x)$ --- найти значение по ключу $x$ (или $\mathrm{null}$, если $x$ нет).
	\end{itemize}
	Здесь $U$ --- \emph{универсум} всех возможных ключей (большое множество, например $\mathbb{Z}$).
\end{Definition}

\begin{Remark}{}{}
	Реализация на основе АВЛ-дерева поддерживает все операции за $O(\log n)$, где $n$ --- текущее количество элементов. Хеш-таблица реализует их в среднем за $O(1)$.
\end{Remark}

\subsubsection*{Хеширование произвольных ключей}

Если $U$ --- не целые числа, используется один из двух подходов.
\begin{enumerate}
	\item \textbf{Уникальные объекты} ($x = y \Leftrightarrow$ это один и тот же объект в памяти): каждому объекту сопоставляется его адрес в памяти $\in \mathbb{N}$ (в C++ это буквально указатель).
	\item \textbf{Неуникальные объекты} (строки, кортежи, векторы): сопоставляем им \emph{хеш} (например, хеш Карпа--Рабина). Коллизия возникает с малой вероятностью.
\end{enumerate}

\subsection{Хеш-таблица с открытой адресацией}

\begin{Definition}{}{}
	Хранится таблица $T[0..m-1]$, где $n < m$, $n$ --- текущее число ключей. \emph{Хеш-функция} $h\colon U \to \{0, \ldots, m-1\}$ определяет, в какую ячейку помещается ключ $x \in U$: он помещается в $T[h(x)]$.
\end{Definition}

\begin{Definition}{}{}
	\emph{Коллизия} --- ситуация, когда $h(x) = h(y)$ при $x \neq y$.
\end{Definition}

\begin{Statement}{}{}
	Если $|U| > m$, то коллизии неизбежны при любой хеш-функции $h$.
\end{Statement}

\subsubsection*{Разрешение коллизий линейным зондированием}

При \emph{линейном зондировании} ячейки таблицы просматриваются последовательно начиная с $h(x)$, по кругу. Используются два специальных значения: $\varnothing$ (пустая ячейка) и $\dagger$ (tombstone --- удалённая ячейка).

\begin{figure}[H]
	\centering
	\begin{tikzpicture}[scale=0.85, >=Stealth, font=\small]
		\def\R{2.2}
		\def\N{18}
		\foreach \i in {0,...,17} {
			\pgfmathsetmacro{\angle}{90 - \i * 360/\N}
			\pgfmathsetmacro{\ax}{\R * cos(\angle)}
			\pgfmathsetmacro{\ay}{\R * sin(\angle)}
			\node[draw, circle, inner sep=2pt, minimum size=10pt, fill=white] at (\ax, \ay) {};
		}
		\foreach \i in {2,3,4,5,6} {
			\pgfmathsetmacro{\angle}{90 - \i * 360/\N}
			\pgfmathsetmacro{\ax}{\R * cos(\angle)}
			\pgfmathsetmacro{\ay}{\R * sin(\angle)}
			\node[draw, circle, inner sep=2pt, minimum size=10pt, fill=gray!60] at (\ax, \ay) {};
		}
		\pgfmathsetmacro{\hangle}{90 - 4 * 360/\N}
		\pgfmathsetmacro{\hx}{\R * cos(\hangle)}
		\pgfmathsetmacro{\hy}{\R * sin(\hangle)}
		\draw[->, thick] (\hx + 0.9, \hy + 0.9) -- (\hx + 0.18, \hy + 0.18);
		\node at (\hx + 1.2, \hy + 1.05) {$h(x)$};
		\node at (4.2, -0.3) {серия};
		\draw[->, black!70] (3.5, -0.2) -- (2.6, -0.35);
	\end{tikzpicture}
	\caption{Хеш-таблица $T[0..m-1]$ как кольцевой буфер. Серые ячейки образуют \emph{серию} --- максимальную цепочку подряд занятых ячеек. Стрелка отмечает позицию $h(x)$, с которой начинается зондирование; время операции пропорционально длине хвоста серии от $h(x)$ до её конца.}
\end{figure}

\subsubsection*{Вставка}
\begin{algorithmic}[1]
	\Function{insert}{$x, v$}
	\State $t = h(x)$;\quad $t_{\dagger} = -1$
	\While{$T[t] \neq \varnothing$ \textbf{and} $(T[t] = \dagger$ \textbf{or} $T[t].\mathit{key} \neq x)$}
	\If{$T[t] = \dagger$ \textbf{and} $t_{\dagger} = -1$}
	\State $t_{\dagger} = t$
	\EndIf
	\State $t = (t + 1) \bmod m$
	\EndWhile
	\If{$T[t].\mathit{key} \neq x$}
	\If{$t_{\dagger} \neq -1$}
	\State $t = t_{\dagger}$
	\EndIf
	\State $T[t].\mathit{key} = x$;\quad $n = n + 1$
	\EndIf
	\State $T[t].\mathit{value} = v$
	\If{$2n > m$}
	\State Создать новую пустую таблицу $T'$ размера $2m$
	\State Перехешировать все элементы из $T$ в $T'$ с учётом нового размера
	\State $T = T'$;\quad $m = 2m$
	\EndIf
	\EndFunction
\end{algorithmic}

\subsubsection*{Удаление}

\begin{algorithmic}[1]
	\Function{delete}{$x$}
	\State $t = h(x)$
	\While{$T[t] \neq \varnothing$ \textbf{and} $T[t].\mathit{key} \neq x$}
	\State $t = (t + 1) \bmod m$
	\EndWhile
	\If{$T[t] \neq \varnothing$}
	\State $T[t] = \dagger$ \Comment{$n$ не меняется}
	\EndIf
	\EndFunction
\end{algorithmic}

\subsection{Амортизированное и ожидаемое время работы}

\begin{Definition}{Ожидаемое время работы}{}
	Пусть $t(a)$ --- (случайное) время работы на входе $a$. \emph{Ожидаемое время работы} алгоритма на входах длины $n$:
	\[
	T(n) = \max_{a \in I_n} \mathbf{E}\bigl[t(a)\bigr],
	\]
	где $I_n$ --- множество всех входов длины $n$.
\end{Definition}

\begin{Definition}{Амортизированное время работы}{}
	Говорят, что \emph{амортизированное время работы} операций структуры данных равно $O(f(n))$, если любая последовательность из $n$ операций, применённых на изначально пустой структуре данных, выполняется за время $O(n\, f(n))$.
\end{Definition}

\begin{Statement}{}{}
	Ожидаемое амортизированное время работы операций в хеш-таблице с открытой адресацией равно $O(1)$.
\end{Statement}

\begin{proof}[Доказательство (набросок)]
	Из-за линейности математического ожидания достаточно показать, что суммарное ожидаемое количество шагов при выполнении $n$ операций равно $O(n)$. Перестройка таблицы (при $2n > m$) происходит лишь после того, как число элементов удваивается; амортизированная стоимость перестройки не превышает $O(1)$ на операцию вставки (по аналогии с динамическим массивом). Остаётся оценить ожидаемое число шагов одной операции без перестройки --- этим занимается следующая теорема.
\end{proof}

\subsection{Идеализированное предположение и анализ}

\begin{Definition}{Идеализированное предположение}{}
	Хеш-функция $h$ генерируется случайно и равновероятно из множества \emph{всех} функций $U \to \{0, \ldots, m-1\}$, т.е. для каждого $x \in U$ значение $h(x)$ выбирается независимо и равновероятно из $\{0, \ldots, m-1\}$.
\end{Definition}

\begin{Theorem}{}{}
	При идеализированном предположении ожидаемое время работы каждой операции равно $O(1)$ (если не происходит перестройка таблицы).
\end{Theorem}

\begin{proof}
	Будем считать, что все ключи различны и $m = 2n$. Время работы операции с ключом $x$ пропорционально длине <<хвоста>> серии, начиная с позиции $t = h(x)$. Обозначим:
	\begin{itemize}
		\item \emph{серия} --- максимальная цепочка подряд занятых ячеек $t', t'+1, \ldots, t'+k-1$, где ячейки $t'-1$ и $t'+k$ пусты (нумерация по модулю $m$);
		\item $q_k$ --- вероятность того, что в таблице существует серия длины $k$, начинающаяся в фиксированной позиции.
	\end{itemize}
	
	\textbf{Ожидаемое время.} Пусть $p_\ell$ --- вероятность встретить в позиции $t = h(x)$ <<хвост>> длины $\ell$ некоторой серии длины $k \geq \ell$. Если серия имеет длину $k \geq \ell$, то для хвоста длины $\ell$ в позиции $t$ она обязана начинаться в позиции $t - (k-\ell)$, поэтому
	\[
	p_\ell = \sum_{k=\ell}^{n} q_k.
	\]
	Тогда ожидаемое время равно $O\!\left(\sum_{\ell=0}^{n} p_\ell \cdot \ell\right)$. Преобразуем:
	\[
	\sum_{\ell=0}^{n} p_\ell \cdot \ell
	= \sum_{\ell=0}^{n} \sum_{k=\ell}^{n} q_k \cdot \ell
	= \sum_{k=0}^{n} q_k \cdot (0 + 1 + \cdots + k)
	= \sum_{k=0}^{n} q_k \cdot \Theta(k^2).
	\]
	
	\textbf{Оценка $q_k$.} Если в таблице есть серия в позициях $t', t'+1, \ldots, t'+k-1$, то ровно $k$ ключей отображены в эти $k$ ячеек, а остальные $n-k$ --- вне них. При $m = 2n$:
	\[
	q_k \leq \binom{n}{k}
	\left(\frac{k}{2n}\right)^{k}
	\left(1 - \frac{k}{2n}\right)^{n-k}.
	\]
	Оценим каждый множитель сверху, используя неравенства
	\[
	\binom{n}{k} \leq \frac{n^n}{k^k(n-k)^{n-k}},
	\qquad
	1+x \leq e^x.
	\]
	Первое неравенство следует из бинома Ньютона: $n^n = \bigl(k + (n-k)\bigr)^n = \sum_{i=0}^{n} \binom{n}{i} k^i (n-k)^{n-i}$. Все слагаемые этой суммы неотрицательны, поэтому, оставив лишь слагаемое с $i = k$, получаем $n^n \ge \binom{n}{k} k^k (n-k)^{n-k}$. Тогда
	\begin{align*}
		q_k
		&\leq
		\frac{n^n}{k^k(n-k)^{n-k}}
		\cdot
		\frac{k^k}{(2n)^k}
		\cdot
		\left(\frac{2n-k}{2n}\right)^{n-k}
		=
		\frac{1}{2^k}
		\cdot
		\left(\frac{2n-k}{2(n-k)}\right)^{n-k} \\
		&=
		\frac{1}{2^k}
		\cdot
		\left(1 + \frac{k}{2(n-k)}\right)^{n-k}
		\leq
		\frac{1}{2^k}
		\cdot
		e^{\,k(n-k)/[2(n-k)]}
		=
		\frac{1}{2^k} \cdot e^{k/2}
		=
		\left(\frac{\sqrt{e}}{2}\right)^k.
	\end{align*}
	Поскольку $\tfrac{\sqrt{e}}{2} < 1$, получаем $q_k \leq e^{-\Omega(k)}$. Следовательно,
	\[
	\sum_{k=0}^{n} q_k \cdot \Theta(k^2)
	\leq \Theta\!\left(\sum_{k=0}^{\infty} k^2 \cdot \left(\tfrac{\sqrt{e}}{2}\right)^k\right) = O(1).
	\]
	Таким образом, ожидаемое время операции равно $O(1)$.
\end{proof}